\subsection{A quadratic register model for history naturality}
\label{subsec:quadratic-register-history}

The relative-divisor theorem starts with a polynomial already supplied.  The
following restricted model runs in the other direction: a declared algebraic
register and a projective arithmetic operator force the polynomial by
elimination.  It is deliberately smaller than
Problem~\ref{prob:history-to-divisor-functor}.  In particular, the divisor below
records the possible \emph{endpoint values} of a register; it is not, by itself,
the divisor of histories that reach zero, and its two sheets are not two
inequivalent marked histories.

Put
\[
 k=\Q(\sqrt2),\qquad R=k[t,t^{-1}],\qquad B=\Spec R,
\]
and supply the quadratic register cover
\begin{equation}
 \pi:\widetilde B=\Spec k[u,u^{-1}]\longrightarrow B,
 \qquad t=u^2.
 \label{eq:quadratic-register-cover}
\end{equation}
This cover is extra typed input: extracting a square root is not one of the four
elementary contexts of Paper~I.  Let \(\gamma\in\operatorname{Hist}_4\) be a
marked history in the distinguished Hecke sublanguage and set
\(h=\rho(\gamma)\in\Gamma_4\).  Choose a determinant-one representative
\[
 M_h=\begin{pmatrix}a&b\\c&d\end{pmatrix}
 \in\operatorname{SL}_2(\Z[\sqrt2]).
\]
The endpoint part of the computation graph is the \(B\)-morphism
\begin{equation}
 \iota_h:\widetilde B\longrightarrow\mathbb P^1_B,
 \qquad
 u\longmapsto[au+b:cu+d].
 \label{eq:register-endpoint-map}
\end{equation}
We call the triple \((\widetilde B,u,\iota_h)\) the
\emph{register-decorated endpoint correspondence}, and denote its
scheme-theoretic image by \(D_h\).

\begin{theorem}[Forced projective quadratic]
\label{thm:forced-projective-quadratic}
With the notation above, the following statements hold.
\begin{enumerate}
\item The endpoint divisor is cut out in \(\mathbb P^1_B\), with homogeneous
coordinates \([Z:W]\), by the binary quadratic
\begin{equation}
 \mathcal P_h(Z,W;t)
  =(dZ-bW)^2-t(aW-cZ)^2.
 \label{eq:forced-homogeneous-quadratic}
\end{equation}
It is independent of the scalar representative of \(h\), up to multiplication
by a unit of \(R\).
\item The morphism \(\iota_h\) identifies \(\widetilde B\) with \(D_h\).
Consequently \(D_h\to B\) is an integral finite \(\acute{e}\)tale cover of
degree two, hence finite flat.
\item On the chart \(W=1\), its equation is
\begin{equation}
 P_h(z,t)=(dz-b)^2-t(a-cz)^2.
 \label{eq:forced-affine-quadratic}
\end{equation}
As a polynomial in \(z\) over \(k(t)\), it is quadratic and irreducible, and
\begin{equation}
 \operatorname{disc}_z(P_h)=4t(ad-bc)^2=4t.
 \label{eq:forced-quadratic-discriminant}
\end{equation}
Its leading coefficient is \(d^2-tc^2\), which need not be a unit of \(R\).
Thus the affine equation must not be inserted directly into the monic affine
relative-divisor theorem.
\item After the nonconstant Kummer base change
\(\widetilde B\to B\), the divisor splits into two disjoint sections:
\begin{align}
 \mathcal P_h
  ={}&\bigl(dZ-bW-u(aW-cZ)\bigr)\notag\\
    &\cdot\bigl(dZ-bW+u(aW-cZ)\bigr),
 \label{eq:forced-quadratic-splitting}
\end{align}
whose values are respectively \(h(u)\) and \(h(-u)\).
Extension of the constant field from \(k\) to \(\overline k\) does not split
the generic polynomial: \(t\) is still nonsquare in \(\overline k(t)\).
\item If \(g\in\Gamma_4\), then
\begin{equation}
 D_{g\circ h}=g_{B*}D_h.
 \label{eq:register-divisor-equivariance}
\end{equation}
More explicitly, for
\(M_g=\left(\begin{smallmatrix}p&q\\r&s\end{smallmatrix}\right)\),
\begin{equation}
 \mathcal P_{g\circ h}(Z,W;t)
 =\mathcal P_h(sZ-qW,pW-rZ;t).
 \label{eq:forced-quadratic-composition}
\end{equation}
Hence, for chronological concatenation---first \(\gamma\), then \(\delta\)---
\begin{equation}
 D_{\gamma\delta}=\rho(\delta)_{B*}D_\gamma,
 \qquad
 \rho(\gamma\delta)=\rho(\delta)\rho(\gamma).
 \label{eq:history-divisor-chronology}
\end{equation}
\end{enumerate}
\end{theorem}

\begin{proof}
Write \(\Delta=ad-bc\).  At the projective point
\([Z:W]=[au+b:cu+d]\), the two rows of the adjugate matrix give
\[
 dZ-bW=\Delta u,
 \qquad
 aW-cZ=\Delta.
\]
Eliminating \(u\) from these equations and \(u^2=t\) gives
\eqref{eq:forced-homogeneous-quadratic}.  Scaling \(M_h\) by
\(q\in k^\times\) scales \(\mathcal P_h\) by \(q^2\), so its Cartier divisor
is unchanged.

For the identity operator, the divisor is
\[
 D_1=V(Z^2-tW^2).
\]
It misses \(W=0\), and on \(W=1\) its coordinate ring is
\[
 R[z]/(z^2-t)\cong k[u,u^{-1}].
\]
It is therefore integral and finite free of rank two over \(R\).  Since
\(4t\in R^\times\), it is also \(\acute{e}\)tale.  The automorphism \(h_B\) of
\(\mathbb P^1_B\) carries \(D_1\) to \(D_h\), proving the first two claims.

Expanding \eqref{eq:forced-affine-quadratic} gives coefficients
\[
 A=d^2-tc^2,\qquad B=2(tac-bd),\qquad C=b^2-ta^2.
\]
The leading coefficient is nonzero in \(k(t)\), since otherwise the
transcendence of \(t\) would force \(c=d=0\).  Direct calculation yields
\(B^2-4AC=4t\Delta^2\).  The polynomial \(z^2-t\) is irreducible over
\(k(t)\), and an invertible projective change of coordinates preserves
irreducibility, proving the generic assertion.  Factoring after \(t=u^2\)
gives \eqref{eq:forced-quadratic-splitting}.  The two factors are disjoint
because \(2u\) is a unit.  The same nonsquare argument over
\(\overline k(t)\) proves the constant-field qualification.

Finally, \(D_h=h_B(D_1)\), so \(D_{g\circ h}=g_B(D_h)\).  Substituting the
inverse coordinates
\([sZ-qW:-rZ+pW]\) gives
\eqref{eq:forced-quadratic-composition}.  The chronological rule is precisely
the convention of Paper~I.  These are elementary instances of finite-flat and
finite-\(\acute{e}\)tale base change
\cite{Hartshorne1977AlgebraicGeometry,GrothendieckRaynaud1971SGA1}.
\end{proof}

The geometric generic fiber has two points over
\(\overline{k(t)}\), but the total cover remains integral after extending only
the constants to \(\overline k\).  Equation
\eqref{eq:forced-quadratic-splitting} is instead a splitting after the
nonconstant register cover itself.

\begin{remark}[The affine pole is not a collision]
\label{rem:quadratic-affine-pole}
The projective endpoint is finite on a chosen register branch precisely when
\(cu+d\ne0\).  Requiring both conjugate branches to be finite removes
\[
 V\bigl(\mathrm N_{k(u)/k(t)}(cu+d)\bigr)
 =V(d^2-c^2t)
\]
from \(B\).  If \(c,d\ne0\), then at \(t=(d/c)^2\) one endpoint moves to
infinity and the affine polynomial loses degree.  Its discriminant remains
nonzero, so no two projective roots collide.  The homogeneous divisor is still
finite \(\acute{e}\)tale there.

If instead one compactifies the base to \(\overline B=\Spec k[t]\), the same
projective equation remains integral and finite flat, but it is ramified at
\(t=0\).  Its fiber there is a nonreduced double point, and the ramification is
tame of index two.  Thus the \(\mathbb G_m\)-model has Frobenius but no inertia,
whereas this boundary compactification supplies the single Kummer inertia point.
\end{remark}

\subsection{Time-tagged traces, zero hits, and ordinary poles}
\label{subsec:tagged-history-trace}

The endpoint divisor forgets the intermediate computation.  There is, however,
a restricted decorated object which is still forced by the same marked history.
For
\(\gamma=(g_1,\ldots,g_n)\), write
\[
 h_0=1,
 \qquad
 h_j=\rho(g_1,\ldots,g_j)
      =g_j\circ\cdots\circ g_1.
\]
Define the \emph{time-tagged register trace}
\begin{equation}
 \mathcal T_\gamma
 =\bigsqcup_{j=0}^n
   \bigl(D_{h_j}\times\{j\}\bigr),
 \label{eq:time-tagged-register-trace}
\end{equation}
together with the transition label \(g_{j+1}\) from time \(j\) to time
\(j+1\).  The tags are part of the object; forgetting them merely superposes
endpoint divisors and destroys chronology.

For ordinary arithmetic, define first a divisor on the projective input line.
Every letter \(\mathsf j[x]=-1/x\) requests division by the current value.  Put
\begin{equation}
 \Pi_\gamma
 =\sum_{\{j:\,g_{j+1}=\mathsf j\}} h_j^*[0]
 \quad\text{on }\mathbb P^1_k,
 \label{eq:universal-history-pole-divisor}
\end{equation}
Let \(\Pi_{\gamma,B}=\Pi_\gamma\times_k B\) on \(\mathbb P^1_B\), and let
\(\Pi_\gamma^{(u)}=\iota_1^*\Pi_{\gamma,B}\) be its pullback to the register
cover.  Multiplicity records repeated requests at the same input.

\begin{proposition}[Restricted trace composition and pole cocycle]
\label{prop:trace-pole-cocycle}
For \(q=4\) histories the following hold.
\begin{enumerate}
\item Chronological concatenation appends the trace of \(\delta\), evaluated
from the terminal state of \(\gamma\), to the tagged trace of \(\gamma\).
The transition labels and time reindexing determine this operation uniquely.
\item The ordinary pole divisors satisfy the exact cocycle identity
\begin{equation}
 \Pi_{\gamma\delta}
 =\Pi_\gamma+\rho(\gamma)^*\Pi_\delta.
 \label{eq:history-pole-cocycle}
\end{equation}
The branchwise ordinary-admissible locus is
\(\widetilde B\setminus\operatorname{Supp}(\Pi_\gamma^{(u)})\).
\item If
\(M_{h_j}=\left(\begin{smallmatrix}a_j&b_j\\c_j&d_j\end{smallmatrix}\right)\),
then the parameters at which time \(j\) reaches zero are the intersection of
\(D_{h_j}\) with the zero section \([0:1]\), hence are cut out on \(B\) by
\begin{equation}
 b_j^2-t a_j^2=0.
 \label{eq:stage-zero-hit-locus}
\end{equation}
On \(B=\mathbb G_{m,k}\) this locus is empty when \(a_jb_j=0\), and otherwise
is the single divisor \(t=(b_j/a_j)^2\).
\end{enumerate}
\end{proposition}

\begin{proof}
The prefix operators of \(\gamma\delta\) are first those of \(\gamma\), and
then the maps \(q_j\circ\rho(\gamma)\), where \(q_j\) runs through the prefixes
of \(\delta\).  This proves the trace statement.  A pole occurring in the
second block has input divisor
\(\rho(\gamma)^*(q_j^*[0])\); summing the first and second blocks gives
\eqref{eq:history-pole-cocycle}.  Since translations have no finite pole and
every inversion pole is listed, the complement of the pulled-back support is
exactly the ordinary-admissible locus for a chosen register branch.  Finally,
substitution of \([Z:W]=[0:1]\) in
\eqref{eq:forced-homogeneous-quadratic} gives
\eqref{eq:stage-zero-hit-locus}.
\end{proof}

The proposition supplies a genuine composition law for a \emph{restricted,
decorated correspondence}: the quadratic seed cover is fixed in advance, every
time is tagged, and every transition and pole request is retained.  It is not a
functor from arbitrary marked histories to prime divisors.  In particular,
\(\mathcal T_\gamma\) is generally a disconnected multisection, while
\(D_{\rho(\gamma)}\) is only its terminal prime divisor.

\begin{example}[A neutral operator with nonneutral trace]
\label{ex:omega-four-register-trace}
Let \(\lambda=\sqrt2\), and let
\[
 \omega_4=(\mathsf t_+,\mathsf j,
            \mathsf t_+,\mathsf j,
            \mathsf t_+,\mathsf j,
            \mathsf t_+,\mathsf j).
\]
With the chronological convention of Paper~I,
\(\rho(\omega_4)=(JT)^4=1\).  Hence \(D_{\omega_4}=D_1\), and even the
register-decorated terminal map is the same as for the empty history.  The four
universal inversion requests occur at
\[
 x=-\sqrt2,\qquad x=-1/\sqrt2,\qquad x=0,\qquad x=\infty.
\]
The register \(u\in\mathbb G_m\) omits the last two, so
\begin{equation}
 \Pi_{\omega_4}^{(u)}=[u=-\sqrt2]+[u=-1/\sqrt2],
 \qquad
 \Pi_{\varnothing}^{(u)}=0.
 \label{eq:omega-four-register-poles}
\end{equation}
The maximal base open on which both conjugate register branches are ordinarily
admissible therefore removes \(t=2\) and \(t=1/2\).  The tagged traces are also
different: one has nine time slices and the other only its initial slice.  Thus
endpoint equality does not imply equality of ordinary domains or computation
traces.
\end{example}

\subsection{What the divisor retains, and what it necessarily forgets}
\label{subsec:register-divisor-kernels}

\begin{proposition}[Exact collapse levels in the quadratic model]
\label{prop:quadratic-register-collapse}
Let \(h_1,h_2\in\Gamma_4\).
\begin{enumerate}
\item As abstract covers of \(B\), all \(D_h\to B\) are isomorphic to the same
Kummer cover \(\widetilde B\to B\).  Their monodromy, ramification behavior, and
Frobenius cycle types therefore do not depend on \(h\).  The operator enters only
through the embedding into the fixed output line \(\mathbb P^1_B\).
\item Let \(\epsilon(z)=-z\).  Then
\begin{equation}
 D_{h_1}=D_{h_2}
 \quad\Longleftrightarrow\quad
 h_2^{-1}h_1\in\{1,\epsilon\}
 \quad\text{in }\operatorname{PGL}_2(k).
 \label{eq:embedded-divisor-kernel}
\end{equation}
Since \(\epsilon\) sends the upper half-plane to the lower half-plane,
\(\Gamma_4\cap\{1,\epsilon\}=\{1\}\).  Thus the embedded divisor recovers the
operator inside \(\Gamma_4\), although the abstract cover does not.
\item Equality of register-decorated endpoint maps
\(\iota_{h_1}=\iota_{h_2}\) implies \(h_1=h_2\).  Nevertheless the assignment
\(\gamma\mapsto(\widetilde B,u,\iota_{\rho(\gamma)})\) factors through
\(\rho\).  Every pair of distinct histories with the same projective operator,
including \(\omega_4\) and the empty word, has the same terminal
correspondence.
\item At a single chosen geometric register value \(u_0\), endpoint equality is
coarser:
\begin{equation}
 h_1(u_0)=h_2(u_0)
 \quad\Longleftrightarrow\quad
 h_2^{-1}h_1\in\operatorname{Stab}_{\Gamma_4}(u_0).
 \label{eq:closed-endpoint-kernel}
\end{equation}
By contrast, equality at the generic register \(u\) forces operator equality.
\end{enumerate}
\end{proposition}

\begin{proof}
The first assertion follows from \(D_h=h_B(D_1)\).  A constant projective
automorphism preserving \(D_1\) restricts to a \(B\)-automorphism of the quadratic
cover.  That cover has only the identity and the deck involution \(u\mapsto-u\).
Because \(u\) is transcendental over \(k\), the corresponding rational maps are
respectively \(z\mapsto z\) and \(z\mapsto-z\).  This proves
\eqref{eq:embedded-divisor-kernel}.  The orientation statement gives its
intersection with \(\Gamma_4\).  If \(h_1(u)=h_2(u)\) in \(k(u)\), two
fractional linear maps agree as rational functions and hence define the same
projective operator.  The final assertion is the usual orbit--stabilizer
criterion.
\end{proof}

This proposition is the plain-divisor no-go in its precise form.  Forgetting the
embedding collapses all operators.  Keeping the embedding can recover the
restricted \(\Gamma_4\) operator, but not a marked word in its evaluation fiber.
Keeping the register distinguishes the two Kummer lifts; keeping the tagged
prefix trace and pole divisor additionally retains intermediate states and
ordinary domain failures.  None of these statements proves that the trace is a
complete invariant of literal or marked-history equality.

\subsection{Frobenius in the quadratic model}
\label{subsec:quadratic-register-frobenius}

Let \(\mathcal O_k=\Z[\sqrt2]\).  Because the standard generators of
\(\Gamma_4\) lie in \(\operatorname{SL}_2(\mathcal O_k)\), the preceding
construction has an integral model over
\(\mathcal O_k[1/2,t,t^{-1}]\).

\begin{proposition}[Quadratic Frobenius cycle]
\label{prop:quadratic-register-frobenius}
Let \(b\) be a closed point of this integral base, let
\(\kappa(b)=\F_q\), and write \(t_b\in\F_q^\times\) for the value of the
register parameter.  On the two geometric endpoint values of \(D_h\), either
arithmetic or geometric Frobenius has cycle type
\begin{equation}
 \begin{cases}
 (1)(1),&t_b\text{ is a square in }\F_q,\\
 (12),&t_b\text{ is a nonsquare in }\F_q.
 \end{cases}
 \label{eq:quadratic-frobenius-cycle}
\end{equation}
This cycle type is independent of \(h\); the embedding merely sends
\(\{\pm\sqrt{t_b}\}\) to \(\{h(\pm\sqrt{t_b})\}\).
\end{proposition}

\begin{proof}
For a root \(r^2=t_b\) in \(\overline{\F}_q\), one has
\(r^q=t_b^{(q-1)/2}r\).  Thus Frobenius fixes the two signs when the quadratic
character is \(+1\) and interchanges them when it is \(-1\).  In a two-element
permutation group arithmetic and geometric Frobenius have the same cycle type.
The coefficients of \(h\) lie in the residue field, so Frobenius commutes with
the projective coordinate change \cite{Neukirch1999AlgebraicNumberTheory}.
\end{proof}

For example, at the prime
\((7,\sqrt2-3)\) the residue field is \(\F_7\).  The specialization \(t=2\)
splits because \(2=3^2\), whereas \(t=3\) gives a transposition.  This arithmetic
cycle belongs to the supplied register cover, not to a new history invariant.

\subsection{A quartic test separating constant and geometric components}
\label{subsec:quartic-register-test}

The quadratic model is arithmetically prime and remains integral after every
extension of constants.  A second forced register graph exhibits a genuinely
different layer.  Retain \(k=\Q(\sqrt2)\), put \(\eta=3\), and supply registers
\begin{equation}
 v^2=\eta,
 \qquad
 u^2=t+v.
 \label{eq:quartic-register-tower}
\end{equation}
The constant \(\eta\) is not a square in \(k\).  Eliminating \(v\) forces
\begin{equation}
 Q(u,t)=(u^2-t)^2-\eta
       =u^4-2tu^2+t^2-\eta.
 \label{eq:forced-register-quartic}
\end{equation}

\begin{theorem}[Arithmetic prime with two constant-geometric components]
\label{thm:forced-register-quartic}
Let \(E\subset\mathbb P^1_B\) be cut out by
\begin{equation}
 \mathcal Q(U,V;t)=(U^2-tV^2)^2-\eta V^4.
 \label{eq:forced-binary-quartic}
\end{equation}
Then:
\begin{enumerate}
\item \(E\to B\) is an integral finite flat cover of degree four, and
\(Q(u,t)\) is irreducible over \(k(t)\).
\item Over \(\overline k(t)\), it has exactly two irreducible quadratic factors
\begin{equation}
 Q=(u^2-(t+\sqrt\eta))(u^2-(t-\sqrt\eta)).
 \label{eq:quartic-constant-factorization}
\end{equation}
Each factor remains irreducible over \(\overline k(t)\).  It splits into two
sections only after the further nonconstant Kummer base change adjoining
\(\sqrt{t\pm\sqrt\eta}\).
\item Its discriminant in \(u\) is
\begin{equation}
 \operatorname{disc}_u Q
 =256\eta^2(t^2-\eta).
 \label{eq:quartic-register-discriminant}
\end{equation}
Thus the cover is finite \(\acute{e}\)tale away from \(t^2=\eta\).  Constant
Galois interchanges the two factors in
\eqref{eq:quartic-constant-factorization}, while analytic monodromy around
\(t=-\sqrt\eta\) or \(t=+\sqrt\eta\) acts within the corresponding quadratic
cover by \(u\mapsto-u\).
\item For \(h\) represented by \(M_h\), put
\(L_1=dZ-bW\) and \(L_0=aW-cZ\).  The projective endpoint divisor of
\(h(u)\) is the binary quartic
\begin{equation}
 \mathcal Q_h(Z,W;t)
 =\bigl(L_1^2-tL_0^2\bigr)^2-\eta L_0^4.
 \label{eq:quartic-projective-pushforward}
\end{equation}
It equals \(h_{B*}E\), so the same composition rule as
\eqref{eq:register-divisor-equivariance} holds.
\end{enumerate}
\end{theorem}

\begin{proof}
The form \(\mathcal Q\) has no point over \(V=0\), and its affine coordinate
ring is the free rank-four \(R\)-module defined by the monic polynomial \(Q\).
In the tower
\[
 k(t)\subset k(\sqrt\eta)(t)
 \subset k(\sqrt\eta)(t)\bigl(\sqrt{t+\sqrt\eta}\bigr),
\]
the first extension has degree two because \(\eta\) is nonsquare in \(k\), and
the second has degree two because \(t+\sqrt\eta\) has a simple zero and hence is
not a square.  The total degree is four, so \(Q\) is the minimal polynomial of
\(u\) over \(k(t)\).  This also proves integrality.

After extending constants, \(\sqrt\eta\) is available and gives
\eqref{eq:quartic-constant-factorization}; the simple-zero argument proves that
both quadratics remain irreducible.  For a polynomial \(u^4+pu^2+r\), the
discriminant is \(16r(p^2-4r)^2\).  Substitution of
\(p=-2t\) and \(r=t^2-\eta\) gives
\eqref{eq:quartic-register-discriminant}.  The Galois and local monodromy
statements follow directly from the two displayed square roots.  Finally,
\([L_1:L_0]=[u:1]\) on the graph of \(h(u)\); homogenized elimination gives
\eqref{eq:quartic-projective-pushforward}.
\end{proof}

There is also an explicit good-reduction dictionary.  Work over
\(\mathcal O_k[1/6,t,t^{-1}]\) and away from \(t^2=\eta\).  At a closed point
with residue field \(\F_q\), let \(\chi\) denote its quadratic character.
If \(\chi(\eta)=+1\), choose \(s\in\F_q\) with \(s^2=\eta\).  Frobenius
preserves the two quadratic components, and on the component
\(u^2=t_b\pm s\) it is the identity or a transposition according as
\(\chi(t_b\pm s)=+1\) or \(-1\).  If \(\chi(\eta)=-1\), Frobenius exchanges
the two constant components.  Choosing \(s\in\F_{q^2}\) with \(s^2=\eta\),
its square acts within \(u^2=t_b+s\) by the character
\begin{equation}
 \chi_{\F_{q^2}}(t_b+s)=\chi_{\F_q}(t_b^2-\eta).
 \label{eq:quartic-frobenius-inner-character}
\end{equation}
Consequently Frobenius has cycle type \((2)(2)\) when the right-hand side is
\(+1\), and a four-cycle when it is \(-1\).  Projective pushforward by \(h\)
does not change these cycle types.

The quartic example supplies the algebraic part of the test requested in the
arithmetic-versus-geometric zero-path problem: arithmetic irreducibility,
constant-field components, local sheet monodromy, discriminant, and Frobenius
are all explicit.  It does not yet supply an AEG assignment or prove that these
registers arise canonically from unrestricted expression histories.

\subsection{Boundary of the restricted result}
\label{subsec:history-divisor-restricted-boundary}

The two constructions establish an exact but narrow chain
\[
 \boxed{
 \begin{gathered}
 \text{supplied algebraic register}
 \longrightarrow
 \text{marked }q=4\text{ projective computation}\\
 \mathord\downarrow\\[-2pt]
 \text{tagged trace and pole divisor}
 \longrightarrow
 \text{projective endpoint divisor}
 \end{gathered}}
\]
The first three arrows retain the declared register, time, transitions, and
ordinary-domain audit.  Forgetting to the final divisor retains at most the
embedded projective operator and its arithmetic cover; forgetting the embedding
retains only the cover.  Thus this section gives a restricted
\(\Gamma_4\)-equivariant assignment, not the general history-to-divisor functor.

\begin{problem}[Descent of trace data without tautological history labels]
\label{prob:trace-decoration-descent}
Find a category of typed algebraic registers and marked arithmetic histories in
which the tagged trace and pole cocycle descend under a declared nontrivial
history equivalence, while retaining enough information to distinguish ordinary
domain failures and zero-hit stages.  Determine whether any finite prime divisor,
normalization, or stateful local system can carry this information without simply
reattaching the complete word as a label.  Then compare the resulting construction
with the automorphic \(q=4\) zero network and with parameter concatenation.
\end{problem}

Until this descent problem is solved, the present results should be read as a
proved minimal naturality model and a proved obstruction: the terminal endpoint
divisor constructed here, because it factors through \(\rho\), cannot be
history-faithful, even though a computation-graph decoration can retain strictly
more than its terminal operator.
